(1) 切向动量定理为 \(mL\ddot\theta=-mg\sin\theta-bL\dot\theta\)。在 \(|\theta|\ll1\) 下除以 \(mL\) 得 \(\ddot\theta+2\gamma\dot\theta+\omega_0^2\theta=0\)，其中 \(\gamma=b/(2m)\)、\(\omega_0^2=g/L\)；阻尼系数单位为每秒。
(2) 当 \(\gamma<\omega_0\) 时特征根为 \(-\gamma\pm i\sqrt{\omega_0^2-\gamma^2}\)。解为 \(\theta=Ae^{-\gamma t}\cos(\omega_dt+\phi)\)，\(\omega_d=\sqrt{\omega_0^2-\gamma^2}\)；\(A,\phi\) 由两个初始条件确定，\(\gamma\ll\omega_0\) 时缓慢变化包络近似有效。
(3) 机械能 \(E=mL^2\dot\theta^2/2+mgL\theta^2/2\)。将运动方程乘以 \(mL^2\dot\theta\) 得 \(\dot E=-bL^2\dot\theta^2\)；平均一个弱阻尼周期给 \(\overline E\propto e^{-2\gamma t}\)。定义 \(Q=2\pi\overline E/|\Delta E_{\rm cycle}|\)，得 \(Q\simeq\omega_0/(2\gamma)=m\omega_0/b\)。
(4) 振幅包络的比值 \(A(t)/A(0)=e^{-\gamma t}\)，因此半振幅时间 \(t_{1/2}=\ln2/\gamma\)。临界阻尼时 \(\gamma=\omega_0\)，解为 \((C_1+C_2t)e^{-\gamma t}\)，没有真实振荡周期，不能将上述振荡品质因数外推到临界点。
